\section{A shared exhaustive signed cover}\label{sec:shared-signed-cover}
The signed reduction permits an arithmetic exclusion to be shared among
all unsigned rings compatible with a signed table. The following
construction makes that shared search exhaustive.

\begin{proposition}[A shared exhaustive signed cover]
\label{prop:global-signed-cover}
Fix $M\geq1$. Consider the integral self-dual rank-four signed tables
with the prescribed unit and reciprocity, parameter tuple
$(a,b,c,d,e,f,g,h,i,j)$, and all ten coefficients of absolute value
at most $M$. The six associativity equations define a finite set
preserved by the group of nonunit basis permutations and basis-sign
changes, of order $48$.

Every orbit has a representative satisfying
\[
 b,c,d\geq0,\qquad
 b=\min\{|b|,|c|,|d|,|f|,|h|,|i|\},
 \qquad e\geq0\ \text{if }b=c=0.
\]
Enumerating all such representatives and retaining only the
lexicographically least preferred tuple in each of their $48$
transforms gives exactly one table per orbit.

For every self-dual unsigned fusion ring $R$ of multiplicity at
most $M$, all its possible sphericalization quotients are covered
by the transforms of these representatives that satisfy its signed
coefficient bounds, parities, and adjoint-parity conditions.
A representative failing the formal-codegree $d$-number criterion
or character cyclotomicity can be discarded for every $R$
simultaneously.
\end{proposition}
\begin{proof}
The permutation group is transitive on the six mixed positions
$b,c,d,f,h,i$. Place a position of minimum absolute value at $b$.
Under a basis-sign change $(s_1,s_2,s_3)$, the signs multiplying
$b,c,d$ are respectively $s_2,s_3,s_1$. Hence these three entries
can be made nonnegative independently, without changing their
absolute values. If $b=c=0$, the remaining choices of $s_2,s_3$
also make $e$ nonnegative while preserving $d\geq0$.
Thus the preferred set meets every orbit. Its least element in
each finite orbit is unique; checking all $48$ transforms selects
precisely that element.

For completeness, the preferred set is exhausted by the following
disjoint branches of the associativity equations. If $b>0$, then
$c,d\geq b$ and $|f|,|h|,|i|\geq b$; the two-coefficient
Diophantine relation \eqref{eq:census-linear}, with the polynomial
coefficients from Appendix~\ref{sec:census}, holds for signed
parameters as well. Exhaust the bounded choices of $b,c,d,e,f$
and solve that relation for all bounded integer pairs $a,i$ by
exact greatest-common-divisor progressions. If one coefficient
vanishes, exhaust its free variable; if both vanish, retain the
entire bounded product precisely when the constant term vanishes.
The formulas in that appendix then determine $h,g,j$. Impose the
signed bounds and the preferred-representative conditions on all
coordinates, and check the original six equations.
If $b=0<c$ and $e\ne0$, the equations determine
\begin{align*}
 h&=\frac{d^2+e^2-ad-1}{c},&
 i&=\frac{e(d+f-a)}c,\\
 j&=\frac{c^2+e^2+f^2-af-1}{c},&
 g&=\frac{h(d-f)+ei-cd}{e}.
\end{align*}
The divisibility for $i$ gives
$a\equiv d+f\pmod{c/\gcd(c,e)}$; all bounded choices in this
progression are checked. If $b=0<c$, $e=0$ and $d>0$, then
\[
 \gcd(c,d)=1,\quad i=0,\quad
 h=\frac{d^2-ad-1}{c}\ne0,\quad f=d-\frac{cd}{h},
\]
and $j$ is given by the preceding formula, while $g$ is enumerated.
The remaining case $b=0<c$, $e=d=0$ forces
$(c,h,f,i,j)=(1,-1,0,0,0)$ and leaves $a,g$ free.

If $b=c=0<e$, then $a=d+f$ and $df=e^2-1$;
the equations determine $g,j$ from $d,e,f,h,i$, with the case
$d=0$ retained. Finally, $b=c=e=0$ forces $d=1,a=0$ and
$f\in\{1,-1\}$. For $f=1$ the remaining equation is
\[
 gi+hj=h^2+i^2-2,
\]
including its $h=0$ branch. For $f=-1$ one has $h=i=0$
and $g,j$ are free. Coordinates whose signs are not fixed by the
preferred convention retain both signs and the original size
bounds, and every emitted tuple is checked
against all six equations. This proves exhaustive enumeration
of the preferred set, including the zero denominators.

Any sphericalization quotient for $R$ lies in the bounded signed
set by Proposition~\ref{prop:signed}. One of the $48$ transforms
of its global representative is the quotient in the basis of $R$,
so testing all compatible transforms loses no categorical
possibility. Formal codegrees are unchanged by permutations and
basis signs, and character values change only by permutation and
sign. Their $d$-number and cyclotomicity conditions are therefore
orbit-invariant. Both conditions are necessary for a signed
sphericalization quotient by Corollary~\ref{cor:joint-signed-conditions},
which proves the common exclusions.
\end{proof}

\begin{remark}[Certificate reuse]
The global signed census is regenerated exhaustively during the
verification workflow; its completeness is not inferred from a
cached list. Universally excluded orbits receive a checked common
arithmetic rejection. Every other orbit is expanded into all
distinct transforms compatible with each unsigned ring, and every
such candidate receives the corresponding signed-branch check.
Sharing these computations changes their allocation among records,
not the necessary criteria or the scope of the classification.
\end{remark}
